\section{总结与延伸}

本讲把 security leadership 从“CSO 是否足够勇敢”重写为一套可验证的组织设计：

\begin{enumerate}
  \item \textbf{安全是 operating system}：people、process、platform 和 governance 共同决定结果；
  \item \textbf{公共规则需要反馈}：企业应把真实系统 evidence 带入政策，同时接受公共 accountability；
  \item \textbf{VDP 管报告渠道}：scope、acknowledgement、修复和协调披露构成信任合同；
  \item \textbf{Bounty 管可选激励}：奖励、NDA 或 payment 不能事后授权有害行为；
  \item \textbf{Incident command 管跨职能执行}：IC 同步工作流，但不垄断技术、法律和经营判断；
  \item \textbf{Evidence 与 decision log 管可复核性}：事实、分析、决定和修正必须分层记录；
  \item \textbf{Materiality 不等于 severity}：对适用上市公司，Form 8-K 的四工作日通常从重大性确定后计算；
  \item \textbf{案件状态必须更新}：课堂的 pending appeal 已被 2025 年第九巡回法院裁判和 2026 年 cert denial 取代；
  \item \textbf{Accountability 共享但不稀释}：CSO、GC、CEO 和 board 各有明确 deliverable 与 decision right；
  \item \textbf{Resilience 与 public service 是长期能力}：恢复人员、保留制度记忆并改善公共技术治理。
\end{enumerate}

\begin{importantbox}{Security Governance 作业}
为一个虚构 SaaS 公司设计端到端 playbook：公开 VDP 与 safe-harbor 边界；定义 VDP-to-incident triggers；画出 IC、technical、legal/comms、business RACI；提供 evidence chain 与 decision-log 模板；列出 customer、regulator、law-enforcement 和 public-market disclosure workstream；最后设计六个 tabletop inject。每个关键决定必须标明 owner、evidence、deadline 和 review condition。
\end{importantbox}

\begin{knowledgebox}{最终自检：五个不能混淆}
不要把 vulnerability 与 incident 混淆，不要把 severity 与 legal materiality 混淆，不要把 payment 与 authorization 混淆，不要把 speaker account 与 adjudicated outcome 混淆，也不要把 shared responsibility 与 diffuse accountability 混淆。只要这五组边界清楚，组织就更可能在高压中保持事实和责任完整。
\end{knowledgebox}

\subsection{拓展阅读}

下面的材料不是同一层级：CISA、DOJ 与 NIST 主要帮助设计 VDP 和 incident-response process；SEC 资料解释特定公开市场披露框架；法院意见和 docket 记录案件程序结果。阅读时应先确定自己要回答的是“如何建设控制”“适用什么义务”还是“法院对既有记录作了什么判断”，避免用工程指南代替法律分析，也避免用单一案件代替普遍工程原则。

{\small
\begin{itemize}[nosep,leftmargin=*]
  \item \href{https://www.cisa.gov/vulnerability-disclosure-policy-template}{\textbf{CISA Vulnerability Disclosure Policy Template}}：scope、reporting、expectations 与 coordinated disclosure 示例。
  \item \href{https://www.cisa.gov/news-events/directives/bod-20-01-develop-and-publish-vulnerability-disclosure-policy}{\textbf{CISA BOD 20-01}}：联邦机构 VDP 的政策背景。
  \item \href{https://www.justice.gov/criminal/criminal-ccips/page/file/983996/dl}{\textbf{DOJ Framework for a Vulnerability Disclosure Program}}：授权、good faith 与执法考量框架。
  \item \href{https://csrc.nist.gov/pubs/sp/800/61/r3/final}{\textbf{NIST SP 800-61 Rev. 3}}：将 incident response 融入 CSF 2.0 风险管理活动。
  \item \href{https://www.sec.gov/files/33-11216-fact-sheet.pdf}{\textbf{SEC Cybersecurity Disclosure Final Rule Fact Sheet}}：material incident 与 Form 8-K 时点概览。
  \item \href{https://cdn.ca9.uscourts.gov/datastore/opinions/2025/11/12/23-927.pdf}{\textbf{United States v. Sullivan, No. 23-927}}：第九巡回法院修订意见。
  \item \href{https://www.supremecourt.gov/docket/docketfiles/html/public/25-1082.html}{\textbf{Supreme Court Docket 25-1082}}：certiorari petition 与 2026 年 6 月 29 日拒绝受理记录。
\end{itemize}
}

\begin{knowledgebox}{建议阅读顺序：从可执行控制到裁判边界}
先用 CISA template 检查公开 policy 是否给出 scope、承诺和联系渠道；再用 DOJ framework 审视 good-faith 与授权边界；随后用 NIST SP 800-61 Rev. 3 把报告升级到组织级 incident response；若组织是适用的美国上市公司，再阅读 SEC fact sheet 与 final rule 原文；最后对照 Ninth Circuit opinion 和 Supreme Court docket，确认课堂案件的程序状态。每一步都记录来源日期和适用范围，因为政策、标准与案件状态会继续变化。
\end{knowledgebox}

\begin{importantbox}{Source triangulation 练习}
选择一个假想 incident，分别写出：技术团队从 NIST 得到的运行建议、counsel 需要核验的通知与披露义务、管理层需要作出的经营决定，以及法院材料\textbf{不能}替你回答的问题。若四栏出现同一句结论，通常说明证据层级仍被混在一起。
\end{importantbox}

\end{document}
